\begin{definition}[Real RegSeqRat StreamName dependency triple]
\label{def:real-regseq-streamname-dependency-triple}
A $\RealUp$ RegSeqRat StreamName dependency triple is a finite
$\mathsf{BHist}$ packet
$$
(S,Q,D,R,H,C,P,N)
$$
whose visible rows have the following shape:
\begin{enumerate}
\item $S$ records the inspected $\StreamNameUp$ finite windows and their
  transported stream-name classifier rows;
\item $Q$ records the $\RegSeqRatUp$ readback surface over the same scheduled
  windows;
\item $D$ records the $\DyadicRatCoreUp$ observation ledger read pointwise
  through those windows;
\item $R$ records the $\RealUp$ seal row and its finite endpoint readback;
\item $H$ records componentwise $\hsame$ comparisons for the window,
  readback, dyadic, and seal rows;
\item $C$ records the displayed $\Cont$ route from each public seal read to
  the preceding window, readback, dyadic, and seal coordinates;
\item $P$ records the $\Pkg$ provenance for the four sibling surfaces;
\item $N$ is the local $\NameCert$ row for this dependency packet.
\end{enumerate}
The classifier compares dependency triples only through these displayed rows.
It has no coordinate for a quotient real, selected host limit, host equality,
ambient completeness principle, arithmetic package, order package, or
metric-completeness row.
\end{definition}

\begin{theorem}[Real RegSeqRat StreamName seal-route exhaustion]
\label{thm:real-regseq-streamname-seal-route-exhaustion}
Every accepted $\RealUp$ seal read from a dependency triple
$(S,Q,D,R,H,C,P,N)$ factors through the $\StreamNameUp$ finite window row
$S$, then the $\RegSeqRatUp$ readback row $Q$, then the
$\DyadicRatCoreUp$ ledger row $D$, and only then the $\RealUp$ seal row $R$.
The factorization is transported by the componentwise $\hsame$ row $H$, read
through the displayed $\Cont$ route $C$, and packaged by $P,N$.

Consequently the accepted seal route is exactly the sibling route through
\autoref{ch:concrete-instances-streamname-namecert},
\autoref{ch:concrete-instances-regseqrat-namecert}, and
\autoref{ch:concrete-instances-dyadicratcore-namecert} before the
$\RealUp$ seal is consumed. It exports no quotient real, selected host limit,
host equality, ambient completeness coordinate, arithmetic package, order
package, or metric-completeness row.
\end{theorem}
\begin{proof}
Begin with an accepted seal read from $(S,Q,D,R,H,C,P,N)$. The definition of
the dependency triple exposes only the finite window row $S$, regular
rational readback $Q$, dyadic observation ledger $D$, seal row $R$, and the
transport, route, package, and local naming rows $H,C,P,N$.

Use \autoref{thm:realup-regseqrat-streamname-seal-handoff}. That handoff
identifies the $\RealUp$ seal-facing package as the stream-name classifier
row, transported unary classifier row, finite-window ledger, endpoint
readback, and seal provenance.

Now read the upstream sibling surfaces. The StreamName, RegSeqRat, and
DyadicRatCore chapters named in
\autoref{ch:concrete-instances-streamname-namecert},
\autoref{ch:concrete-instances-regseqrat-namecert}, and
\autoref{ch:concrete-instances-dyadicratcore-namecert} provide exactly the
window, readback, and dyadic ledger rows used by the handoff. The displayed
$H,C,P,N$ rows keep those coordinates in one $\mathsf{BHist}$ packet.

Since the packet has no remaining coordinate after $R$, a public seal read
cannot branch to a quotient real, selected host limit, host equality, ambient
completeness principle, arithmetic or order package, or metric-completeness
row. Thus every accepted read is exhausted by the stated route.
\end{proof}
